% A LaTeX (non-official) template for ISAE projects reports
% Copyright (C) 2014 Damien Roque
% Version: 0.2
% Author: Damien Roque <damien.roque_AT_isae.fr>

\documentclass[10pt,a4paper]{article}
\usepackage[utf8]{inputenc}
\usepackage[T1]{fontenc}
\usepackage[frenchb]{babel} % If you write in French
%\usepackage[english]{babel} % If you write in English
\usepackage{a4wide}
\usepackage{graphicx}
\graphicspath{{images/}}
\usepackage{subfig}
\usepackage{tikz}
\usetikzlibrary{shapes,arrows}
\usepackage{pgfplots}
\pgfplotsset{compat=newest}
\pgfplotsset{plot coordinates/math parser=false}
\newlength\figureheight
\newlength\figurewidth
\pgfkeys{/pgf/number format/.cd,
set decimal separator={,\!},
1000 sep={\,},
}
\usepackage{ifthen}
\usepackage{ifpdf}
\ifpdf
\usepackage[pdftex]{hyperref}
\else
\usepackage{hyperref}
\fi
\usepackage{color}
\hypersetup{%
colorlinks=true,
linkcolor=black,
citecolor=black,
urlcolor=black}

\renewcommand{\baselinestretch}{1.05}
\setlength{\headheight}{15pt}

\let\headruleORIG\headrule
\usepackage{colortbl}
\arrayrulecolor{black}

\makeatletter
\def\@textbottom{\vskip \z@ \@plus 1pt}
\let\@texttop\relax
\makeatother

\makeatletter
\def\cleardoublepage{\clearpage\if@twoside \ifodd\c@page\else%
  \hbox{}%
  \thispagestyle{empty}%
  \newpage%
  \if@twocolumn\hbox{}\newpage\fi\fi\fi}
\makeatother

\usepackage{amsthm}
\usepackage{amssymb,amsmath,bbm}
\usepackage{array}
\usepackage{bm}
\usepackage{multirow}
\usepackage[footnote]{acronym}

\newcommand*{\SET}[1]  {\ensuremath{\mathbf{#1}}}
\newcommand*{\VEC}[1]  {\ensuremath{\boldsymbol{#1}}}
\newcommand*{\FAM}[1]  {\ensuremath{\boldsymbol{#1}}}
\newcommand*{\MAT}[1]  {\ensuremath{\boldsymbol{#1}}}
\newcommand*{\OP}[1]  {\ensuremath{\mathrm{#1}}}
\newcommand*{\NORM}[1]  {\ensuremath{\left\|#1\right\|}}
\newcommand*{\DPR}[2]  {\ensuremath{\left \langle #1,#2 \right \rangle}}
\newcommand*{\calbf}[1]  {\ensuremath{\boldsymbol{\mathcal{#1}}}}
\newcommand*{\shift}[1]  {\ensuremath{\boldsymbol{#1}}}

\newcommand{\eqdef}{\stackrel{\mathrm{def}}{=}}
\newcommand{\argmax}{\operatornamewithlimits{argmax}}
\newcommand{\argmin}{\operatornamewithlimits{argmin}}
\newcommand{\ud}{\, \mathrm{d}}
\newcommand{\vect}{\text{Vect}}
\newcommand{\sinc}{\ensuremath{\mathrm{sinc}}}
\newcommand{\esp}{\ensuremath{\mathbb{E}}}
\newcommand{\hilbert}{\ensuremath{\mathcal{H}}}
\newcommand{\fourier}{\ensuremath{\mathcal{F}}}
\newcommand{\sgn}{\text{sgn}}
\newcommand{\intTT}{\int_{-T}^{T}}
\newcommand{\intT}{\int_{-\frac{T}{2}}^{\frac{T}{2}}}
\newcommand{\intinf}{\int_{-\infty}^{+\infty}}
\newcommand{\Sh}{\ensuremath{\boldsymbol{S}}}
\newcommand{\C}{\mathbb{C}}
\newcommand{\R}{\mathbb{R}}
\newcommand{\Z}{\mathbb{Z}}
\newcommand{\N}{\mathbb{N}}
\newcommand{\K}{\mathbb{K}}
\newcommand{\reel}{\mathcal{R}}
\newcommand{\imag}{\mathcal{I}}
\newcommand{\cmnr}{c_{m,n}^\reel}
\newcommand{\cmni}{c_{m,n}^\imag}
\newcommand{\cnr}{c_{n}^\reel}
\newcommand{\cni}{c_{n}^\imag}
\newcommand{\tproto}{g}
\newcommand{\rproto}{\check{g}}
\newcommand{\LR}{\mathcal{L}_2(\SET{R})}
\newcommand{\LZ}{\ell_2(\SET{Z})}
\newcommand{\LZI}[1]{\ell_2(\SET{#1})}
\newcommand{\LZZ}{\ell_2(\SET{Z}^2)}
\newcommand{\diag}{\operatorname{diag}}
\newcommand{\noise}{z}
\newcommand{\Noise}{Z}
\newcommand{\filtnoise}{\zeta}
\newcommand{\tp}{g}
\newcommand{\rp}{\check{g}}
\newcommand{\TP}{G}
\newcommand{\RP}{\check{G}}
\newcommand{\dmin}{d_{\mathrm{min}}}
\newcommand{\Dmin}{D_{\mathrm{min}}}
\newcommand{\Image}{\ensuremath{\text{Im}}}
\newcommand{\Span}{\ensuremath{\text{Span}}}


\theoremstyle{definition}
\newtheorem{exemple}{Exemple}[section]
\newtheorem{remarque}{Remarque}[section]

\theoremstyle{plain}
\newtheorem{theoreme}{Théorème}[section]
\newtheorem{result}{Résultat}[section]
\newtheorem{definition}{Définition}[section]
\newtheorem{propriete}{Propriété}[section]
\newtheorem{lemme}{Lemme}[section]
\newtheorem{proposition}{Proposition}[section]

\theoremstyle{remark}



\parskip=5pt
%\sloppy

\begin{document}

%%%%%%%%%%%%%%%%%%
%%% First page %%%
%%%%%%%%%%%%%%%%%%

\begin{titlepage}
\begin{center}

{\large Génie mathématique - 4ème année}\\[0.5cm]

{\large Rapport de stage}\\[0.5cm]

% Title
\rule{\linewidth}{0.5mm} \\[0.4cm]
{ \huge \bfseries Étude du comportement des méthodes BFGS et L-BFGS pour résoudre un sous-problème de région de confiance \\[0.4cm] }
\rule{\linewidth}{0.5mm} \\[1.5cm]

% Author and supervisor
\noindent
\begin{minipage}{0.4\textwidth}
  \begin{flushleft} \large
    \emph{Auteur :}\\
    Johann \textsc{Bourhis}\\
    % M. Prénom \textsc{Nom}\\
    % M\up{me} Prénom \textsc{Nom}\\
    % M. Prénom \textsc{Nom}
  \end{flushleft}
\end{minipage}%
\begin{minipage}{0.4\textwidth}
  \begin{flushright} \large
    \emph{Encadrants :} \\
    Jean-Pierre \textsc{Dussault}\\
    \& Dominique \textsc{Orban}\\
    %Dr.~Prénom \textsc{Nom}
  \end{flushright}
\end{minipage}

\vspace{0.3cm}

\includegraphics[width=0.8\textwidth]{logo}\\[1cm]

\includegraphics[width=0.7\textwidth]{Polytechnique}\\[1cm]

\vfill

% Bottom of the page
{\today}

\end{center}
\end{titlepage}

%%%%%%%%%%%%%%%%%%%%%%%%%%%%%
%%% Non-significant pages %%%
%%%%%%%%%%%%%%%%%%%%%%%%%%%%%


%\section*{Remerciements}




\clearpage
\tableofcontents

%\clearpage
%\listoffigures

\section{Introduction}
%\addcontentsline{toc}{section}{Introduction}
%\markboth{Introduction}{Introduction}
\label{sec:introduction}



%%% Local Variables: 
%%% mode: latex
%%% TeX-master: "isae-report-template"
%%% End: 


%\clearpage
%\section*{Notations}

%\\
%\begin{itemize}
%    \item[$\bullet$]    coucou
%\end{itemize}



%%%%%%%%%%%%%%%%%%%%%%%%%%%%%%%%%%%%%%%%%%%%
%%% Content of the report and references %%%
%%%%%%%%%%%%%%%%%%%%%%%%%%%%%%%%%%%%%%%%%%%%

%\mainmatter

\cleardoublepage
\section{Étude de la méthode BFGS} 

\subsection{Construction de la méthode et comportement sur une quadratique convexe}

\subsection{Résultats de convergence et relations avec la méthode du gradient conjugué}

\subsection{Comportement des deux méthodes dans le cas de problèmes non-convexes}

\section{Utilisation de la méthode BFGS avec une contrainte de région de confiance}

\subsection{Implémentation de l'algorithme de région de confiance}


\begin{equation}\tag{SP1}\label{subprob1}
    \min\limits_{\alpha \geq 0}f(x + \alpha d)
\end{equation}

\begin{equation}\tag{SP2}\label{subprob2}
        \min\limits_{p \in \mathbb{R}^n} \{f(x + p) \text{ : } \|p\| \leq \Delta \}.
\end{equation}

\begin{equation}\tag{1}
    \Omega(\Delta) := \{x \in \mathbb{R}^n \text{ tel que } \|x\| \leq \Delta\}
\end{equation}

\subsection{Construction d'une recherche linéaire avec contrainte de région de confiance}
Dans cette sous-section, nous cherchons à adapter la méthode BFGS à une contrainte de région de confiance. L'objectif ici est d'appliquer la recherche linéaire avec la mise à jour de BFGS tant que l'on est à l'intérieur de l'ensemble $\Omega(\Delta)$ et de trouver une stratégie pour calculer un itéré $s_k$ dans le cas où l'on sortirait de cette région de confiance.

%%%
\begin{theoreme}
    Soient $d_k = -H_k g_k$ la direction de descente définie par la méthode BFGS et $\alpha_k^* = \frac{-g_k^t d_k}{d_k^t A d_k}$ la longueur de pas optimale du sous problème (\ref{subprob1}). On définit
\begin{equation}\tag{5}\label{vectors}
    p_n = \sum_{k=1}^{n-1} s_k,
\end{equation}
\begin{center}
    où $s_k = \alpha^*_k d_k$.
\end{center}
Alors la suite des itérés ${\|p_n\|}$ est croissante et la fonction $f$ décroît le long du chemin formé par les itérés $s_k$.
\end{theoreme}

\noindent
\textit{Preuve.} D'après la propriété de conjugaison nous savons que
\begin{equation}\tag{1}\label{ineq}
    \|p_{n+1}\|_A^2 = \|p_n + s_n\|_A^2 = \|p_n\|_A^2 + \|s_n\|_A^2 > \|p_n\|_A^2.
\end{equation}
Soient $\{v_1, \text{ } v_2, \text{ }..., \text{ } v_n\}$ les vecteurs propres de $A$ associés respectivement aux valeurs propres $\{\lambda_1, \text{ } \lambda_2, \text{ }..., \text{ } \lambda_n\}$. Comme $A$ est inversible, nous pouvons écrire $p_n$ et $p_{n+1}$ comme étant des combinaisons linéaires des $\{v_k\}_{1 \leq k \leq n}$
\begin{equation}\tag{2.a}
    p_{n+1} = \sum_{i=1}^n \beta_i v_i,
\end{equation}
\vspace{-0.2cm}
\begin{equation}\tag{2.b}
    p_n = \sum_{i=1}^n \alpha_i v_i.
\end{equation}
Étant donné que les $\{v_k\}_{1 \leq k \leq n}$ forment une famille orthogonale, il en suit que
\begin{equation}\tag{3.a}\label{orthogo}
    \|p_{n+1}\|_A^2 = \sum_{i=1}^n \beta_i^2 \lambda_i^2 \|v_i\|^2,
\end{equation}
\vspace{-0.2cm}
\begin{equation}\tag{3.b}
    \|p_n\|_A^2 = \sum_{i=1}^n \alpha_i^2 \lambda_i^2 \|v_i\|^2.
\end{equation}
Ainsi (\ref{ineq}) implique que
\begin{equation}\tag{4}
    \|p_{n+1}\|^2 = \sum_{i=1}^n \beta_i^2 \|v_i\|^2 > \sum_{i=1}^n \alpha_i^2 \|v_i\|^2 = \|p_n\|^2.
\end{equation}
Nous avons par ailleurs l'inégalité suivante
\begin{equation}\tag{5}
    \|p_{n+1}\|^2 = \|p_n\|^2 + \|s_n\|^2 + 2p_n^t s_n > \|p_n\|^2 \iff \frac{-p_n^t s_n}{\|s_n\|^2} < \frac{1}{2}.
\end{equation}

Nous venons ainsi de montrer que la suite $\{\|p_n\|\}_{n \geq 1}$ est croissante et donc que les itérés $s_k$ s'éloignent à chaque étape du centre de la région de confiance $\Omega(\Delta)$. Il y a alors deux possiblités d'arrêt de l'algorithme en tenant compte de la contrainte de région de confiance. Soit $f(p_n)$ converge vers le minimum, dans le cas où celui-ci se trouve à l'intérieur de $\Omega(\Delta)$, soit $p_n$ sort de la région de confiance et il faut alors trouver le vecteur $p^*$ s'approchant au mieux de la solution avec contrainte.

%%%%%%%%

Pour un $k \in \mathbb{N}^*$ donné, posons la fonction $h\text{ : }\mathbb{R} \rightarrow \mathbb{R}$
\begin{equation}\tag{1}
    h(\alpha) = f(p_k + \alpha d_k) = \frac{1}{2}(p_k +\alpha d_k)^t A(p_k +\alpha d_k) + b^t(p_k +\alpha d_k)
\end{equation}
Il en suit que
\begin{equation}\tag{2.a}
    h'(\alpha) = d_k^t A(p_k +\alpha d_k) + b^t d_k
\end{equation}
\vspace{-0.5cm}
\begin{equation}\tag{2.b}
    h''(\alpha) = d_k^t A d_k.
\end{equation}
Etant donné que $A$ est SDP, $h"(\alpha) > 0$ et donc $h'(\alpha)$ est croissante. Nous savons donc que $h'(\alpha)$ est négative $\forall \alpha < \alpha^* := \frac{-\nabla f(p_k)^t d_k}{d_k^t A d_k}$.

%%%
\noindent
Or il se trouve que $\alpha^*$ est le pas optimal du sous-problème (\ref{subprob1}). Ainsi, le chemin constitué des itérés $\{d_k\}_{k \in \mathbb{N}^*}$ est de la forme 
\begin{equation}\tag{4}
    \mathcal{C}(k, \alpha) = \alpha_{k-1}^*d_{k-1} + \alpha d_k, \text{ } k \in \mathbb{N}^*, \text{ } \alpha \in [0, \alpha_k^*],
\end{equation}
\begin{center}
    avec $d_0 = 0$, $\alpha_0 = 0$.
\end{center}

Le résultat précédent nous indique donc que la fonction $f$ décroit le long du chemin $\mathcal{C}(k, \alpha)$. Nous pouvons ainsi choisir, s'il existe, le point au croisement entre $\mathcal{C}(k, \alpha)$ et $\Omega(\Delta)$ comme solution approchée du sous-problème (\ref{subprob2}). Nous sommes alors assurés d'avoir fait décroître au mieux $f$ dans la dernière direction $d_k$ donnée, tout en respectant la contrainte de région de confiance.

%%%
\noindent
\begin{result}
    Supposons donc que le minimum atteint par $f$ ne se situe pas dans la région de confiance $\Omega(\Delta)$. Nous reprenons la définition (\ref{vectors}) de $p_n$ et $s_n$. Soit $n$ le premier entier tel que $\|p_{n+1}\| > \Delta$. Nous savons d'après le précédent théorème qu'il existe $\tau^* \in \text{} ]0, 1[$ tel que 
    \begin{equation}\tag{6} \label{norme_Delta}
        \|p_n + \tau^* s_n \| = \Delta.
    \end{equation}
\end{result}
Tout d'abord élevons la norme au carré. L'équation (\ref{norme_Delta}) se résoud alors en trouvant la racine positive de
\begin{equation}\tag{7}
    \tau^2 \|s_n\|^2 + 2 \tau p_n^t s_n + (\|p_n\|^2 - \Delta^2).
\end{equation}
Cette équation possède deux solutions réelles distinctes puisque $(\|p_n\|^2 - \Delta^2) < 0$ par hypothèse. En effet, suivant la direction donnée par $s_n$, il est possible d'atteindre la frontière de $\Omega(\Delta)$ dans le sens positif et dans le sens négatif.
Les deux racines nous sont données par
\begin{equation}\tag{8}
    \tau_{1, 2} = \frac{-p_n^t s_{n} \pm \sqrt{(p_n^t s_{n})^2 + \|s_n\|^2(\Delta^2 - \|p_n\|^2)}}{\|s_n\|^2}.
\end{equation}
Il est clair que $\sqrt{(p_n^t s_{n})^2 + \|s_n\|^2(\Delta^2 - \|p_n\|^2)} > p_n^t s_{n}$. Il existe donc une unique racine positive qui nous est donnée par
\begin{equation}\tag{9}
    \tau^* = \frac{-p_n^t s_{n} + \sqrt{(p_n^t s_{n})^2 + \|s_n\|^2(\Delta^2 - \|p_n\|^2)}}{\|s_n\|^2}.
\end{equation}
Comme voulu, $\tau^* \in \text{ } ]0, 1[$. En effet, posons
\begin{equation}\tag{10}
    h(\tau) = \tau^2 \|s_n\|^2 + 2 \tau p_n^t s_n + \|p_n\|^2.
\end{equation}
Il en suit que
\begin{equation}\tag{11.a}
    h'(\tau) = 2\tau \|s_n\|^2 + 2p_n^ts_n,
\end{equation}
\vspace{-0.5cm}
\begin{equation}\tag{11.b}
    h''(\tau) = 2\|s_n\|^2.
\end{equation}
Comme $h''(\tau) > 0$, nous avons que $h'(\tau)$ est croissante et donc positive à partir de $\hat{\tau} = \frac{-p_n^ts_n}{\|s_n\|^2} \leq \frac{1}{2}$. Puisque $h(1) = \|p_n + s_n\|^2 > \Delta^2 = h(\tau^*)$, nous avons enfin la relation $\hat{\tau} < \tau^* < 1$.

%%%
\noindent
Nous pouvons donc établir un protocole à suivre pour appliquer la méthode BFGS à un sous-problème de région de confiance. Tant que le minimum n'est pas atteint, nous choisissons $s_k$ et $p_{k+1}$ tels qu'indiqués par (\ref{vectors}). Si $\|p_{k+1}\| > \Delta$, nous remplaçons $p_{k+1}$ par $p_k + \tau^* s_k$ et arrêtons la recherche linéaire.

\begin{remarque}
    Rien ne nous assure ici que nous avons bien trouvé le minimum du sous problème (\ref{subprob2}). Mais cela n'a que très peu d'importance. En effet, nous cherchons en réalité à trouver la valeur minimale de $f$ suivant les directions $d_k$ données par la méthode BFGS. Puisque nous sommes assurés que $p_n$ s'éloigne du centre de $\Omega(\Delta)$, nous savons que nous nous sommes rapprochés de son minimum comme désiré.
\end{remarque}


%%%
\subsection{Intérêt de la région de confiance dans le cas non-convexe}

\section{Implémentation de la méthode L-BFGS et comparaison avec les autres méthodes}

\subsection{Une méthode quasi-Newton à mémoire limitée}

\subsection{Comparaison avec les méthodes Quasi-Newton}

\subsection{Comparaison avec la méthode du gradient conjugué}

%%\section*{Conclusion}
%\addcontentsline{toc}{chapter}{Conclusion}
%\markboth{Conclusion}{Conclusion}
%\label{sec:conclusion}
\section{Annexes}





%\appendix

\clearpage

\bibliographystyle{type_bib}
\bibliography{references}
\nocite{*}

\clearpage

%%%%%%%%%%%%%%%%
%%% Abstract %%%
%%%%%%%%%%%%%%%%

\thispagestyle{empty}

\vspace*{\fill}
\noindent\rule[2pt]{\textwidth}{0.5pt}\\
{\textbf{Abstract ---}}


{\textbf{Key words :}}


\noindent\rule[2pt]{\textwidth}{0.5pt}
\begin{center}
  INSA Rennes\\
  Génie Mathématique\\
  Année 2018-2019\\
\end{center}
\vspace*{\fill}

\end{document}